\chapter{Estado modular bicapa y resolvente espectral}
\label{rm:chap:estado-modular-resolvente}

\section{El operador angular como antecedente com\'un}

El operador constitutivo del \cref{rm:ch:vacuum} admite una formulaci\'on
previa en la que intensidad y orientaci\'on ocupan coordenadas distintas de
un mismo objeto. Sea \(R=R^*=R^{-1}\) una involuci\'on de multiplicidades
espectrales uno y uno, y sean
\begin{equation}
 P_+=\frac{I+R}{2},
 \qquad
 P_-=\frac{I-R}{2}
 \label{rm:eq:modular-projectors}
\end{equation}
sus proyectores. Para
\begin{equation}
 x=A_{\rm rad}>0,
 \qquad
 y=C^*_{\rm rad},
 \qquad
 |y|<x,
 \label{rm:eq:modular-domain}
\end{equation}
definimos el generador angular y su exponencial positiva por
\begin{equation}
 K_{\angle}=xI+yR,
 \qquad
 T_{x,y}=\ee^{-K_{\angle}}
 =q_+P_++q_-P_-,
 \label{rm:eq:modular-angular-operator}
\end{equation}
donde
\begin{equation}
 q_+=\ee^{-x-y},
 \qquad
 q_-=\ee^{-x+y}.
 \label{rm:eq:modular-qpm}
\end{equation}
La condici\'on \(|y|<x\) sit\'ua ambos autovalores en \((0,1)\). El
determinante y el cociente espectral separan las dos coordenadas:
\begin{equation}
 D_A:=\det T_{x,y}=q_+q_-=\ee^{-2x},
 \qquad
 \varrho_{\rm or}:=\frac{q_-}{q_+}=\ee^{2y}.
 \label{rm:eq:modular-determinant-ratio}
\end{equation}
El primer car\'acter es par bajo el intercambio de hojas; el segundo registra
su orientaci\'on. La recuperaci\'on es inmediata:
\begin{equation}
 x=-\frac12\log D_A,
 \qquad
 y=\frac12\log\varrho_{\rm or}.
 \label{rm:eq:modular-recover-xy}
\end{equation}

La procedencia causal de \(x\) y \(y\) permanece vinculada al estado
\(\mathrm{APP}\to\TRIT\to\TPK\). En particular,
\begin{equation}
 x=\frac{50\pi}{9}\,\mathfrak A_\beta,
 \qquad
 D_A=\exp\!\left(-\frac{100\pi}{9}\mathfrak A_\beta\right),
 \label{rm:eq:modular-alpha-determinant}
\end{equation}
donde \(\mathfrak A_\beta=\alpha_{\HMT}\) denota el car\'acter terminal del
macroobjeto dodecaf\'asico ya construido. La coordenada \(y\) procede despu\'es
del lector regional y conserva la orientaci\'on que el determinante elimina.

\begin{proposition}[Separaci\'on can\'onica de intensidad y orientaci\'on]
\label{rm:prop:modular-separation}
La aplicaci\'on
\[
 (x,y)\longmapsto(D_A,\varrho_{\rm or})
 =(\ee^{-2x},\ee^{2y})
\]
es un difeomorfismo entre el dominio \(x>0\), \(|y|<x\), y el dominio
\[
 0<D_A<1,
 \qquad
 D_A<\varrho_{\rm or}<D_A^{-1}.
\]
La involuci\'on de hojas act\'ua mediante
\(
(D_A,\varrho_{\rm or})\mapsto(D_A,\varrho_{\rm or}^{-1})
\).
\end{proposition}

\begin{proof}
Las f\'ormulas inversas son \cref{rm:eq:modular-recover-xy}. La desigualdad
\(|y|<x\) equivale a
\(|\log\varrho_{\rm or}|<-\log D_A\), y de ella se obtiene el intervalo
declarado. Sustituir \(y\) por \(-y\) conserva \(D_A\) e invierte el
cociente.
\end{proof}

\section{Normalizaci\'on del estado bicapa}

La traza del operador angular es
\begin{equation}
 \Tr T_{x,y}=2\ee^{-x}\cosh y.
 \label{rm:eq:modular-trace-T}
\end{equation}
Su normalizaci\'on elimina exactamente la intensidad \(x\) y conserva la
orientaci\'on:
\begin{equation}
 \boxed{
 \rho_y:=\frac{T_{x,y}}{\Tr T_{x,y}}
 =\frac{\ee^{-yR}}{2\cosh y}
 =p_+(y)P_++p_-(y)P_- ,}
 \label{rm:eq:modular-rho}
\end{equation}
con
\begin{equation}
 p_+(y)=\frac{\ee^{-y}}{2\cosh y},
 \qquad
 p_-(y)=\frac{\ee^{y}}{2\cosh y},
 \qquad
 p_++p_-=1.
 \label{rm:eq:modular-probabilities}
\end{equation}
La familia \(\rho_y\) es fiel para todo \(y\in\R\). Su observable de
orientaci\'on satisface
\begin{equation}
 \Tr(\rho_yR)=-\tanh y,
 \qquad
 \rho_{-y}=R_{\rm sh}\rho_yR_{\rm sh}^*,
 \label{rm:eq:modular-orientation-expectation}
\end{equation}
donde \(R_{\rm sh}\) representa cualquier unitario que permute
\(P_+\) y \(P_-\). Esta implementaci\'on distingue la involuci\'on
espectral \(R\), que mide el signo de hoja, del unitario \(R_{\rm sh}\),
que intercambia las hojas.

\begin{definition}[Hamiltoniano hologr\'afico modular local]
\label{rm:def:holographic-modular-hamiltonian}
El Hamiltoniano modular del estado bicapa es
\begin{equation}
 \boxed{
 K_y^{\rm mod}:=-\log\rho_y
 =\log(2\cosh y)I+yR.}
 \label{rm:eq:modular-Hamiltonian}
\end{equation}
La componente escalar normaliza el estado y la componente \(yR\) conserva
la memoria orientada.
\end{definition}

El adjetivo \emph{hologr\'afico} registra la procedencia de \(R\),
\(P_\pm\) y \(y\) desde el estado enriquecido y sus dos hojas. La igualdad
\cref{rm:eq:modular-Hamiltonian} es una identidad de c\'alculo funcional; su
promoci\'on a generador de una din\'amica temporal requiere especificar el
flujo f\'isico correspondiente.

\section{Entrop\'ia, paridad y divergencia de orientaci\'on}

\begin{theorem}[Geometr\'ia informacional exacta de las dos hojas]
\label{rm:thm:modular-information-geometry}
La entrop\'ia de von Neumann, la paridad y la entrop\'ia relativa entre las
dos orientaciones satisfacen
\begin{align}
 S(\rho_y)
 &:=-\Tr(\rho_y\log\rho_y)
 =\log(2\cosh y)-y\tanh y,
 \label{rm:eq:modular-entropy}\\
 S(\rho_y)&=S(\rho_{-y}),
 \label{rm:eq:modular-entropy-even}\\
 \boxed{
 D(\rho_y\Vert\rho_{-y})
 =2y\tanh y.}
 \label{rm:eq:modular-relative-entropy}
\end{align}
La divergencia es positiva para \(y\ne0\) y se anula exactamente en la
hoja no orientada \(y=0\).
\end{theorem}

\begin{proof}
De \cref{rm:eq:modular-Hamiltonian} y
\(\Tr(\rho_yR)=-\tanh y\) resulta
\[
 S(\rho_y)=\Tr(\rho_yK_y^{\rm mod})
 =\log(2\cosh y)-y\tanh y.
\]
Ambos t\'erminos son pares. Adem\'as,
\[
 \log\rho_y-\log\rho_{-y}=-2yR,
\]
y, por consiguiente,
\[
 D(\rho_y\Vert\rho_{-y})
 =-2y\Tr(\rho_yR)=2y\tanh y.
\]
El producto \(y\tanh y\) es positivo para \(y\ne0\).
\end{proof}

La entrop\'ia cuantifica la distribuci\'on par de las hojas, mientras la
divergencia cuantifica la separaci\'on informacional de sus orientaciones.
Las derivadas
\begin{equation}
 \frac{\dd}{\dd y}S(\rho_y)=-y\operatorname{sech}^2y,
 \qquad
 \frac{\dd^2}{\dd y^2}\log(2\cosh y)
 =\operatorname{sech}^2y
 \label{rm:eq:modular-entropy-derivatives}
\end{equation}
identifican \(S(\rho_0)=\log2\) como m\'aximo y proporcionan la informaci\'on
de Fisher de la familia exponencial. En el entorno de la hoja sim\'etrica,
\begin{equation}
 D(\rho_y\Vert\rho_{-y})
 =2y^2-\frac23y^4+O(y^6).
 \label{rm:eq:modular-relative-expansion}
\end{equation}
La primera contribuci\'on es cuadr\'atica: una memoria orientada peque\~na
posee un coste informacional de segundo orden.

La distancia de traza entre orientaciones tambi\'en se obtiene de forma
cerrada:
\begin{equation}
 \frac12\lVert\rho_y-\rho_{-y}\rVert_1=|\tanh y|.
 \label{rm:eq:modular-trace-distance}
\end{equation}
Como \(|y|\geq|\tanh y|\), la identidad
\cref{rm:eq:modular-relative-entropy} implica el control
\begin{equation}
 D(\rho_y\Vert\rho_{-y})
 \geq 2\left(\frac12
 \lVert\rho_y-\rho_{-y}\rVert_1\right)^2,
 \label{rm:eq:modular-pinsker-exact-control}
\end{equation}
compatible con la desigualdad de Pinsker. As\'i, la cantidad orientada posee
simult\'aneamente una lectura espectral, entr\'opica y m\'etrica.

\section{Resolvente normalizada del macroobjeto angular}

Sea
\begin{equation}
 \kappa_\beta:=4\pi\mathfrak A_\beta>0.
 \label{rm:eq:modular-kappa-beta}
\end{equation}
Definimos la funci\'on matricial meromorfa
\begin{equation}
 \boxed{
 \mathfrak M_\beta(z)
 :=\kappa_\beta^{-1}(zI-T_{x,y})^{-1},
 \qquad
 z\in\C\setminus\{q_+,q_-\}.}
 \label{rm:eq:modular-resolvent}
\end{equation}
Su descomposici\'on espectral es
\begin{equation}
 \mathfrak M_\beta(z)
 =\frac1{\kappa_\beta}
 \left(\frac{P_+}{z-q_+}+\frac{P_-}{z-q_-}\right).
 \label{rm:eq:modular-resolvent-spectral}
\end{equation}
La resolvente re\'une en un solo objeto la escala de acoplamiento, los polos
angulares, los proyectores de hoja, el determinante y todo el c\'alculo
funcional posterior.

\begin{theorem}[Inyectividad por los dos primeros momentos]
\label{rm:thm:modular-resolvent-injective}
La asignaci\'on
\[
 (\mathfrak A_\beta,T_{x,y})
 \longmapsto\mathfrak M_\beta
\]
es inyectiva para \(\mathfrak A_\beta>0\). En concreto, si
\begin{align}
 B_0&:=\lim_{|z|\to\infty}z\mathfrak M_\beta(z),
 \label{rm:eq:modular-B0}\\
 B_1&:=\lim_{|z|\to\infty}z^2
 \left(\mathfrak M_\beta(z)-z^{-1}B_0\right),
 \label{rm:eq:modular-B1}
\end{align}
entonces
\begin{equation}
 B_0=\kappa_\beta^{-1}I,
 \qquad
 \kappa_\beta=\tau_2(B_0)^{-1},
 \qquad
 T_{x,y}=\kappa_\beta B_1,
 \label{rm:eq:modular-reconstruct-kappa-T}
\end{equation}
donde \(\tau_2=\Tr/2\). Por tanto, la asint\'otica recupera
\(\mathfrak A_\beta=\kappa_\beta/(4\pi)\) y el momento siguiente recupera
el operador angular completo.
\end{theorem}

\begin{proof}
Para \(|z|>\lVert T_{x,y}\rVert\), la serie de Neumann da
\begin{equation}
 \mathfrak M_\beta(z)
 =\frac1{\kappa_\beta}
 \sum_{n\geq0}\frac{T_{x,y}^{n}}{z^{n+1}}.
 \label{rm:eq:modular-Laurent}
\end{equation}
Los coeficientes de \(z^{-1}\) y \(z^{-2}\) son precisamente \(B_0\) y
\(B_1\). Las igualdades \cref{rm:eq:modular-reconstruct-kappa-T} siguen al
tomar la traza normalizada y multiplicar por \(\kappa_\beta\). Dos pares con
la misma resolvente poseen los mismos coeficientes y, por ello, coinciden.
\end{proof}

Si \(q_+\ne q_-\), los datos espectrales se recuperan tambi\'en de los polos
y residuos:
\begin{equation}
 \operatorname*{Res}_{z=q_\pm}\mathfrak M_\beta(z)
 =\kappa_\beta^{-1}P_\pm,
 \qquad
 \det\mathfrak M_\beta(z)
 =\frac{\kappa_\beta^{-2}}
 {(z-q_+)(z-q_-)}.
 \label{rm:eq:modular-resolvent-residues}
\end{equation}
En particular, el denominador escalar contiene
\(q_++q_-\) y \(q_+q_-=D_A\), y
\begin{equation}
 P_+=\frac{T-q_-I}{q_+-q_-},
 \qquad
 P_-=\frac{T-q_+I}{q_--q_+}.
 \label{rm:eq:modular-projector-recovery}
\end{equation}
El caso \(y=0\) produce un \'unico polo simple de multiplicidad espectral
dos: \(T=\ee^{-x}I\). La resolvente conserva entonces la intensidad y el
operador escalar, mientras la elecci\'on de proyectores carece de contenido
intr\'inseco porque la orientaci\'on es exactamente nula.

\begin{corollary}[C\'alculo funcional desde la resolvente]
\label{rm:cor:modular-functional-calculus}
Para toda funci\'on \(f\) holomorfa en un entorno de
\(\{q_+,q_-\}\),
\begin{equation}
 f(T)=\frac{\kappa_\beta}{2\pi\ii}
 \int_\Gamma f(z)\mathfrak M_\beta(z)\,\dd z,
 \label{rm:eq:modular-Dunford}
\end{equation}
donde \(\Gamma\) rodea ambos polos. En particular, la respuesta del vac\'io
\[
 \mathcal V=(I-T^{90})(I-T^{120})^{-1}
\]
se reconstruye de \(\mathfrak M_\beta\) sin incorporar datos escalares
adicionales.
\end{corollary}

\begin{proof}
La identidad es la f\'ormula de Dunford--Riesz aplicada a
\((zI-T)^{-1}=\kappa_\beta\mathfrak M_\beta(z)\). La funci\'on
\(f(z)=(1-z^{90})/(1-z^{120})\) es holomorfa sobre un entorno del espectro,
pues \(0<q_\pm<1\).
\end{proof}

\section{Compatibilidad con el refinamiento non\'adico}

La informaci\'on anterior se prolonga por la torre
\begin{equation}
 \mathcal A_m=M_2(\C)\otimes M_{9^m}(\C),
 \qquad
 \iota_m(a)=a\otimes I_9.
 \label{rm:eq:modular-UHF-tower}
\end{equation}
Sean
\begin{equation}
 T_m=T\otimes I_{9^m},
 \quad
 R_m=R\otimes I_{9^m},
 \quad
 \mathfrak M_{\beta,m}(z)
 =\kappa_\beta^{-1}(zI-T_m)^{-1},
 \label{rm:eq:modular-UHF-objects}
\end{equation}
y sea
\(E_m=\operatorname{id}_{M_2\otimes M_{9^m}}\otimes\tau_9\) la
expectativa tracial desde \(\mathcal A_{m+1}\) hacia \(\mathcal A_m\).

\begin{theorem}[Naturalidad resolvente y modular]
\label{rm:thm:modular-UHF-naturality}
Para todo \(m\geq0\) y todo \(z\) del conjunto resolvente,
\begin{align}
 E_m(T_{m+1})&=T_m,
 \label{rm:eq:modular-UHF-T}\\
 E_m\!\left(\mathfrak M_{\beta,m+1}(z)\right)
 &=\mathfrak M_{\beta,m}(z),
 \label{rm:eq:modular-UHF-resolvent}\\
 E_m\!\left(f(T_{m+1})\right)&=f(T_m)
 \label{rm:eq:modular-UHF-functional}
\end{align}
para toda funci\'on continua \(f\) sobre el espectro. Asimismo, las
funcionales de estado
\begin{equation}
 \omega_{y,m}(a)
 :=(\Tr_2\otimes\tau_{9^m})
 \bigl((\rho_y\otimes I_{9^m})a\bigr)
 \label{rm:eq:modular-UHF-state}
\end{equation}
son compatibles con las inclusiones y definen un estado \(\omega_{y,\infty}\)
sobre el l\'imite UHF.
\end{theorem}

\begin{proof}
Todos los objetos de \cref{rm:eq:modular-UHF-objects} tienen la forma
\(a\otimes I_{9^m}\). La expectativa tracial elimina el \'ultimo factor
identidad y conserva \(a\). La misma observaci\'on prueba
\cref{rm:eq:modular-UHF-functional}; la compatibilidad de los estados se
obtiene evaluando \(\omega_{y,m+1}(a\otimes I_9)\) y usando
\(\tau_9(I_9)=1\).
\end{proof}

La entrop\'ia local relativa a la traza
\(\Tr_2\otimes\tau_{9^m}\) permanece igual a
\(S(\rho_y)\), y la divergencia entre orientaciones permanece igual a
\(2y\tanh y\). El refinamiento a\~nade resoluci\'on non\'adica y conserva el
contenido bicapa. De esta manera, la construcci\'on vive en la torre UHF y
en su cierre tracial, en lugar de depender de una presentaci\'on accidental
por matrices \(2\times2\).

\section{Exactitud de la resolvente modular, obstrucciones de normalización y separación de la lectura termodinámica}

\begin{proofstatusbox}
La descomposici\'on de \(T\), el estado \(\rho_y\), el Hamiltoniano modular,
la entrop\'ia, la divergencia relativa, la inyectividad de la resolvente y
la naturalidad UHF son resultados exactos una vez fijados
\((\mathfrak A_\beta,C^*,R)\). La reuni\'on de esas identidades en la
resolvente \(\mathfrak M_\beta\) constituye una formalizaci\'on nueva de la
arquitectura autoral preexistente. La lectura de
\(D(\rho_y\Vert\rho_{-y})\) como trabajo f\'isico requiere temperatura,
protocolo y din\'amica; hasta fijarlos, su estatuto es el de medida
informacional exacta. Las realizaciones mediante Barbero, CKM o Higgs son
lectoras posteriores y conservan sus respectivos entrelazadores.
\end{proofstatusbox}

\begin{falsifierbox}
El bloque queda refutado si la normalizaci\'on de \(T\) conserva dependencia
de \(x\), si \(K_y^{\rm mod}\) no exponentia a \(\rho_y\), si la entrop\'ia
relativa difiere de \(2y\tanh y\), si dos pares distintos
\((\mathfrak A_\beta,T)\) poseen la misma resolvente, o si alguna expectativa
\(E_m\) altera el c\'alculo funcional. La recuperaci\'on de proyectores en
\(y=0\) constituye adem\'as un control negativo: cualquier selector de hoja
que sobreviva a esa degeneraci\'on habr\'ia sido introducido exteriormente.
\end{falsifierbox}
